\chapter{Introdução}

A ascensão das redes sociais transformou radicalmente a maneira como a informação é consumida e distribuída globalmente. No entanto, essa democratização trouxe consigo o fenômeno da desinformação em escala industrial, caracterizando o que muitos autores chamam de "era da pós-verdade". Nesse cenário, fatos objetivos têm menos influência na formação da opinião pública do que apelos à emoção e crenças pessoais.

\section{Contextualização e o Desafio das IAs Generativas}

Historicamente, a detecção de \textit{fake news} tem sido abordada predominantemente por meio do Processamento de Linguagem Natural (NLP). Esses métodos clássicos buscam padrões linguísticos, inconsistências semânticas ou tons sensacionalistas no corpo do texto para classificar a veracidade de uma notícia. Contudo, essa abordagem encontra-se atualmente saturada e vulnerável. Com a popularização de Inteligências Artificiais generativas de texto avançadas, tornou-se trivial fabricar desinformação com gramática impecável e coesão estrutural, burlando facilmente os detectores baseados exclusivamente em análise textual.

\section{O Problema do "Cold Start" em Novas Redes}

O desafio da moderação de conteúdo torna-se ainda mais crítico em plataformas emergentes ou descentralizadas, como o Bluesky, operado via \textit{AT Protocol}. Diferente de redes maduras como o X (antigo Twitter), que possuem décadas de dados históricos e reputação consolidada de usuários, redes recém-nascidas sofrem com o problema do \textit{Cold Start} (Partida Fria). Quando uma informação falsa começa a circular em uma rede nova, os algoritmos tradicionais não possuem contexto prévio suficiente sobre os emissores para rotular a postagem com precisão e rapidez, permitindo que a desinformação ganhe tração antes de ser mitigada.

\section{Objetivos da Pesquisa}

Diante dessas limitações, este Trabalho de Conclusão de Curso propõe uma mudança de paradigma: focar não apenas no "o que é dito" (texto), mas principalmente em "quem diz" e "como a informação se espalha" (topologia). 

O objetivo principal deste estudo é desenvolver e avaliar um sistema de detecção de \textit{fake news} baseado em Redes Neurais de Grafos (GNNs). Ao invés de analisar publicações isoladas, o modelo avalia o comportamento da rede e as bolhas ideológicas, mapeando os nós (usuários) e as arestas (relações de \textit{follow}, \textit{repost} e \textit{like}) para classificar a árvore de propagação inteira. 

Como objetivos específicos, o trabalho busca:
\begin{itemize}
    \item Extrair e estruturar dados nativos da rede Bluesky via API \textit{atproto}, construindo grafos de interação;
    \item Implementar uma arquitetura de \textit{Graph Convolutional Network} (GCN) baseada no \textit{benchmark} UPFD (\textit{User Preference-aware Fake News Detection});
    \item Comparar o desempenho preditivo de modelos focados puramente em texto com modelos que integram texto e topologia de rede em ambientes de \textit{Cold Start}.
\end{itemize}

\section{Estrutura do Trabalho}

Para atingir os objetivos propostos, este documento está organizado da seguinte forma: o Capítulo 2 apresenta a fundamentação teórica necessária, abordando conceitos de NLP, Teoria de Grafos e \textit{Deep Learning} aplicado a grafos. O Capítulo 3 detalha a metodologia, incluindo a arquitetura do sistema construído, a coleta do \textit{dataset} no Bluesky e a engenharia de dados aplicada. O Capítulo 4 expõe os resultados obtidos, apresentando o duelo entre os modelos e o estudo de ablação, seguido de uma discussão sobre as descobertas. Por fim, o Capítulo 5 traz as considerações finais, limitações encontradas e sugestões para trabalhos futuros.